\section{Билет 4. Начальный план и усложнения ТЗ}

\ticketq{Методы построения начального плана. Усложнения транспортной задачи.}

\paragraph{Что требуется от начального плана.}
Для старта метода потенциалов нужен допустимый \emph{базисный} план закрытой
задачи: $\sum_j x_{ij} = a_i$, $\sum_i x_{ij} = b_j$, $x_{ij} \ge 0$ и ровно
$m+n-1$ заполненных клеток (ранг системы $m+n-1$: одно из $m+n$ уравнений
зависимо). Клеток меньше — план вырожден, добавляют нулевые базисные перевозки.

\paragraph{Метод северо-западного угла.}
Заполнение с клетки $(1,1)$, тарифы не учитываются.
\begin{csteps}
  \item В текущую клетку $(i,j)$ записать $x_{ij} = \min\{a_i, b_j\}$.
  \item $a_i := a_i - x_{ij}$, $b_j := b_j - x_{ij}$.
  \item $a_i = 0$ — к следующей строке; $b_j = 0$ — к следующему столбцу; при
        $a_i = b_j = 0$ закрыть одну строку и один столбец, добавив при
        необходимости нулевую базисную перевозку.
  \item Повторять до распределения всего; проверить число клеток $m+n-1$.
\end{csteps}

\paragraph{Метод минимального элемента.}
То же, но клетка выбирается по тарифу: среди незакрытых берётся $(i,j)$ с
минимальным $c_{ij}$ (при равенстве — самая северо-западная), далее шаги 1--4
дословно те же. Даёт обычно более дешёвый план, так как сначала используются
дешёвые маршруты, но оптимальности не гарантирует.

\paragraph{Оба метода дают базисный план.}
Выкладка общая, различается только правило выбора клетки.
\[
  a_i, b_j \ge 0 \ \Rightarrow\ x_{ij} = \min\{a_i, b_j\} \ge 0 ;
\]
каждый шаг закрывает хотя бы строку или столбец $\Rightarrow$ процесс конечен,
а заполненных клеток не более $m+n-1$ (на последнем шаге закрываются обе);
задача закрытая ($\sum_i a_i = \sum_j b_j$) $\Rightarrow$ исчерпание всех
запасов означает удовлетворение всех потребностей $\Rightarrow$ выполнены
$\sum_j x_{ij} = a_i$, $\sum_i x_{ij} = b_j$, $x_{ij} \ge 0$ — план допустим.

Он к тому же \emph{базисный}: заполненные клетки образуют «ступенчатую»
цепочку (каждый шаг сдвигается на строку вниз или на столбец вправо и назад
не возвращается), поэтому цикла среди них нет, то есть соответствующие столбцы
линейно независимы. $\blacksquare$

\paragraph{Проверка плана.}
$\sum_j x_{ij} = a_i$; $\sum_i x_{ij} = b_j$; $x_{ij} \ge 0$; число базисных
клеток $m+n-1$.

\paragraph{Усложнения транспортной задачи.}
\begin{csteps}
  \item \textbf{Открытая модель} $\sum_i a_i \ne \sum_j b_j$ — приводится к
        закрытой фиктивным участником с нулевыми тарифами:
        \[
          \sum_i a_i > \sum_j b_j :\quad
            b_{n+1} = \sum_i a_i - \sum_j b_j ,\ c_{i,n+1} = 0 ;
        \]
        \[
          \sum_i a_i < \sum_j b_j :\quad
            a_{m+1} = \sum_j b_j - \sum_i a_i ,\ c_{m+1,j} = 0 .
        \]
        Перевозки в фиктивный пункт — это нераспределённый запас (или
        неудовлетворённая потребность); нулевые тарифы не меняют $F$, поэтому
        оптимум исходной задачи сохраняется.
  \item \textbf{Вырожденность:} заполненных клеток меньше $m+n-1$ — добавляют
        нулевые базисные перевозки так, чтобы не образовался цикл.
  \item \textbf{Многопродуктовая модель.} Для продукта $k$ вводятся
        $x_{ij}^{(k)}$:
        $\sum_j x_{ij}^{(k)} = a_i^{(k)}$,
        $\sum_i x_{ij}^{(k)} = b_j^{(k)}$,
        $\sum_k x_{ij}^{(k)} \le d_{ij}$ (пропускная способность маршрута).
  \item \textbf{Запрет поставки} из $A_i$ в $B_j$: клетку исключают либо
        задают $c_{ij} = M$ (большое число). При достаточно большом $M$
        положительная перевозка по этой клетке увеличила бы $F$ сильнее любого
        допустимого варианта, поэтому в оптимальном плане она нулевая — если
        допустимый план без неё вообще существует.
  \item \textbf{Обязательная поставка:} $x_{ij} > 0$, на практике заменяется
        на $x_{ij} \ge K_{ij}$, $K_{ij} > 0$.
  \item \textbf{Не менее $K_{ij}$:} объём выделяют заранее,
        $a_i := a_i - K_{ij}$, $b_j := b_j - K_{ij}$ (сдвиг
        $x_{ij}' = x_{ij} - K_{ij}$) — задача снова обычная.
  \item \textbf{Не более $K_{ij}$:} $x_{ij} \le K_{ij}$ — пропускная
        способность.
  \item \textbf{Сменно-суточный характер:} индекс времени $x_{ij}^{(t)}$;
        учитываются смены, сутки, время доставки, доступность транспорта.
  \item \textbf{Возврат тары, порожний пробег:} обратные перевозки
        $x_{ji}^{\text{back}}$, возможна задача минимизации порожнего пробега.
\end{csteps}

Пункты 1--5 не выводят задачу из класса транспортных: запрет и верхняя граница
— ограничения на отдельные клетки, нижняя снимается сдвигом переменной.
Пункты 6--7 увеличивают размерность и требуют уже общего симплекс-метода, а не
табличной схемы.

\seefull{гл.~2, \S3 и \S8 (численные примеры в \S3)}
